% !TEX program = xelatex
\documentclass[11pt,a4paper]{article}

\usepackage[UTF8]{ctex}
\usepackage[margin=2.2cm]{geometry}
\usepackage{amsmath, amssymb, amsthm}
\usepackage{listings}
\usepackage{xcolor}
\usepackage{tikz}
\usepackage{booktabs}
\usepackage{multirow}
\usepackage{hyperref}
\usetikzlibrary{positioning, arrows.meta, shapes.geometric, fit, calc, decorations.pathreplacing}

\hypersetup{
  colorlinks=true,
  linkcolor=blue!60!black,
  urlcolor=teal!70!black,
}

\lstdefinelanguage{Rust}{
  morekeywords={fn,let,mut,pub,struct,impl,use,mod,self,Self,Option,Some,None,return,if,else,for,while,loop,match,break,continue,as,where,trait,type,enum,move,ref,const,static,unsafe,extern,crate,super,dyn,async,await,box,in},
  sensitive=true,
  morecomment=[l]{//},
  morecomment=[s]{/*}{*/},
  morestring=[b]",
}

\lstset{
  language=Rust,
  basicstyle=\ttfamily\footnotesize,
  keywordstyle=\color{blue!60!black}\bfseries,
  commentstyle=\color{green!40!black}\itshape,
  stringstyle=\color{red!60!black},
  numbers=left,
  numberstyle=\tiny\color{gray},
  numbersep=8pt,
  frame=single,
  framerule=0.4pt,
  rulecolor=\color{gray!40},
  breaklines=true,
  breakatwhitespace=true,
  showstringspaces=false,
  tabsize=4,
  columns=fullflexible,
  keepspaces=true,
}

\title{\textbf{halfedge\,: 拓扑自检模块设计文档} \\ \large \texttt{src/validate.rs}}
\author{halfedge 库}
\date{2026-07-04}

\begin{document}
\maketitle
\tableofcontents

\section{模块定位}

\texttt{validate.rs} 提供 \texttt{validate\_topology}：对 \texttt{MeshStorage} 执行
\textbf{完整}的流形三角曲面不变量校验，输出所有违例的列表（而非首个错误）。

\subsection{与 \texttt{topology\_ops::validate\_mesh} 的区别}

库中存在\textbf{四个语义相关}的验证函数。下表给出明确的取舍标准：

\begin{center}
\renewcommand{\arraystretch}{1.20}
\begin{tabular}{@{}lllll@{}}
  \toprule
  \textbf{函数} & \textbf{错误类型} & \textbf{返回方式} & \textbf{校验项} & \textbf{适用场景} \\
  \midrule
  \texttt{topology\_ops::validate\_mesh}
    & \texttt{TopologyError}
    & 首错返回
    & 仅 twin / next-prev / 三角面环
    & 拓扑操作内部断言（split/flip/collapse） \\
  \texttt{validate::validate\_topology}
    & \texttt{Vec<ValidationError>}
    & 收集全部
    & 全部 8 项
    & 诊断 / 调试：需要所有违例 \\
  \texttt{validate::check\_topology}
    & \texttt{Result<(), Vec>}
    & \texttt{Result} 包装
    & 同 \texttt{validate\_topology}
    & 入口检查：作为前置 / 后置条件，\texttt{?} 传播 \\
  \texttt{validate::validate\_first\_error}
    & \texttt{Option<ValidationError>}
    & 首错返回
    & 同 \texttt{validate\_topology}（但底层复用）
    & 快速失败但需结构化错误 \\
  \bottomrule
\end{tabular}
\end{center}

\paragraph{决策树}
\begin{center}
\begin{tikzpicture}[
  node distance=0.45cm,
  dec/.style={diamond, draw, aspect=2.2, fill=orange!12, inner sep=1pt, font=\small, align=center, minimum width=4.0cm},
  proc/.style={draw, rounded corners=2pt, minimum width=5.0cm, minimum height=0.7cm, align=center, font=\small, fill=green!12},
  op/.style={-Stealth, thick},
  startend/.style={draw, rounded corners=10pt, minimum width=2.6cm, minimum height=0.7cm, font=\small, fill=gray!12}
]
  \node[startend] (s) {开始};
  \node[dec, below=of s] (d1) {在 \texttt{topology\_ops} 内部？};
  \node[proc, below left=0.6cm and 0.2cm of d1] (m) {\texttt{validate\_mesh}（TopologyError）};
  \node[dec, below right=0.6cm and 0.2cm of d1] (d2) {需要看到所有错误？};
  \node[proc, below=0.4cm of d2] (vt) {\texttt{validate\_topology}（Vec）};
  \node[dec, below right=0.6cm and 0.2cm of d2] (d3) {需要 \texttt{?} 传播？};
  \node[proc, below=0.4cm of d3] (ct) {\texttt{check\_topology}（Result）};
  \node[proc, right=0.4cm of d3] (vf) {\texttt{validate\_first\_error}（Option）};
  \draw[op] (s) -- (d1);
  \draw[op] (d1) -| node[above, near start, font=\scriptsize] {是} (m);
  \draw[op] (d1) -| node[above, near start, font=\scriptsize] {否} (d2);
  \draw[op] (d2) -- node[right, font=\scriptsize] {是} (vt);
  \draw[op] (d2) -| node[above, near start, font=\scriptsize] {否} (d3);
  \draw[op] (d3) -- node[right, font=\scriptsize] {是} (ct);
  \draw[op] (d3) -- node[above, font=\scriptsize] {否} (vf);
\end{tikzpicture}
\end{center}

\paragraph{旧表（仅 \texttt{validate\_mesh} vs \texttt{validate\_topology}）}
\begin{center}
\renewcommand{\arraystretch}{1.18}
\begin{tabular}{@{}lll@{}}
  \toprule
  \textbf{特性} & \texttt{validate\_mesh} & \texttt{validate\_topology} \\
  \midrule
  定位     & 拓扑操作前后快速断言 & 完整网格自检 \\
  错误返回 & 首个错误（\texttt{Result}） & 全部错误（\texttt{Vec}） \\
  校验项   & twin / next-prev / 面边界环长度 & 8 项（见下） \\
  退化面   & 不检查 & 检查面积 $< 10^{-12}$ \\
  流形约束 & 不检查 & 检查边与顶点 \\
  入口字段 & 不检查 & 检查 \texttt{Vertex/Face.halfedge} 一致性 \\
  \bottomrule
\end{tabular}
\end{center}

\section{校验项总览}

\begin{center}
\renewcommand{\arraystretch}{1.18}
\begin{tabular}{@{}lll@{}}
  \toprule
  \textbf{\#} & \textbf{校验项} & \textbf{不变量} \\
  \midrule
  1 & 句柄有效性       & 所有 \texttt{vertex/face/twin/next/prev} 引用指向存活元素 \\
  2 & twin 双向匹配    & $A.\mathrm{twin} = B \Rightarrow B.\mathrm{twin} = A$ \\
  3 & next/prev 一致   & $A.\mathrm{next} = B \Leftrightarrow B.\mathrm{prev} = A$ \\
  4 & 面边界环长度     & 每个 \texttt{next} 链闭合，长度 $= 3$ \\
  5 & 入口字段一致     & \texttt{Vertex.halfedge} 的 origin 是该顶点；\texttt{Face.halfedge} 的 face 是该面 \\
  6 & 退化面           & 三角面积 $\ge 10^{-12}$ \\
  7 & 流形边           & 每条无向边至多被 2 个面共享 \\
  8 & 流形顶点         & outgoing 环或闭合（内部），或开链恰 2 个边界端点（边界） \\
  \bottomrule
\end{tabular}
\end{center}

\section{错误类型}

\texttt{ValidationError} 枚举共 14 个变体：

\begin{center}
\renewcommand{\arraystretch}{1.18}
\begin{tabular}{@{}ll@{}}
  \toprule
  \textbf{变体} & \textbf{触发条件} \\
  \midrule
  \texttt{HalfEdgeDanglingVertex} & 半边.\texttt{vertex} 指向已删除顶点 \\
  \texttt{HalfEdgeDanglingFace}   & 半边.\texttt{face} 指向已删除面 \\
  \texttt{HalfEdgeDanglingTwin}   & 半边.\texttt{twin} 指向已删除半边 \\
  \texttt{HalfEdgeDanglingNext}   & 半边.\texttt{next} 指向已删除半边 \\
  \texttt{HalfEdgeDanglingPrev}   & 半边.\texttt{prev} 指向已删除半边 \\
  \texttt{TwinMismatch}           & $A.\mathrm{twin}=B$ 但 $B.\mathrm{twin}\ne A$ \\
  \texttt{SelfLoopHalfEdge}       & $A.\mathrm{vertex} = A.\mathrm{twin}.\mathrm{vertex}$ \\
  \texttt{NextPrevMismatch}       & $A.\mathrm{next}=B$ 但 $B.\mathrm{prev}\ne A$ \\
  \texttt{FaceNotTriangular}      & 面边界环长度 $\ne 3$ \\
  \texttt{VertexHalfEdgeInconsistent} & 顶点入口失效或 origin 不匹配 \\
  \texttt{FaceHalfEdgeInconsistent}   & 面入口失效或 face 不匹配 \\
  \texttt{DegenerateFace}         & 面积 $< 10^{-12}$ \\
  \texttt{NonManifoldEdge}        & 同一无向边被 $> 2$ 个面共享 \\
  \texttt{NonManifoldVertex}      & outgoing 环边界端点数 $\notin \{0, 2\}$ \\
  \bottomrule
\end{tabular}
\end{center}

\section{主校验流程}

\texttt{validate\_topology(mesh)} 依次调用 8 个子校验函数，所有错误汇入
\texttt{Vec<ValidationError>}。便捷包装 \texttt{check\_topology} 返回
\texttt{Result<(), Vec>}，\texttt{validate\_first\_error} 返回
\texttt{Option<ValidationError>}（首个违例或 \texttt{None}）。

\paragraph{\texttt{validate\_first\_error} 的实现取舍}
当前实现为 \texttt{validate\_topology(mesh).into\_iter().next()}，即\textbf{底层
仍跑完全部 8 项校验}，仅在返回时取首元素。这是\textbf{有意的}简化：
\begin{itemize}
  \item 真正的早返回性能由 \texttt{topology\_ops::validate\_mesh} 提供（仅 4 项校验）；
  \item 需要结构化错误（\texttt{ValidationError}）又需要早返回的场景极少，
        现有实现的额外开销仅是\textbf{一次性} \texttt{Vec} 分配；
  \item 若未来确有性能需求，可将 8 个子校验函数重构为
        \texttt{FnMut(ValidationError) -> ControlFlow} 模式，
        \texttt{validate\_topology} 与 \texttt{validate\_first\_error} 共用同一份
        早返回逻辑。
\end{itemize}

\begin{center}
\resizebox{\textwidth}{!}{%
\begin{tikzpicture}[
  node distance=0.45cm,
  proc/.style={draw, rounded corners=2pt, minimum width=10cm, minimum height=0.7cm, align=center, font=\small},
  op/.style={-Stealth, thick},
  startend/.style={draw, rounded corners=10pt, minimum width=2.6cm, minimum height=0.7cm, font=\small, fill=gray!12}
]
  \node[startend] (s) {开始};
  \node[proc, below=of s, fill=blue!8]  (p1) {1. validate\_halfedge\_references：检查 vertex / face / twin / next / prev 存活};
  \node[proc, below=of p1, fill=blue!8] (p2) {2. validate\_twin\_relations：twin 互指 + 无自环};
  \node[proc, below=of p2, fill=blue!8] (p3) {3. validate\_next\_prev：next / prev 反向一致};
  \node[proc, below=of p3, fill=blue!8] (p4) {4. validate\_face\_boundaries：面 next 链闭合，长度 $= 3$};
  \node[proc, below=of p4, fill=blue!8] (p5) {5. validate\_entry\_fields：Vertex / Face 入口字段一致};
  \node[proc, below=of p5, fill=blue!8] (p6) {6. validate\_degenerate\_faces：三角面积 $\ge 10^{-12}$};
  \node[proc, below=of p6, fill=blue!8] (p7) {7. validate\_manifold\_edges：无向边至多 2 面};
  \node[proc, below=of p7, fill=blue!8] (p8) {8. validate\_manifold\_vertices：outgoing 环边界端点数};
  \node[proc, below=of p8, fill=green!12] (pe) {返回 \texttt{Vec<ValidationError>}（空表示通过）};
  \node[startend, below=of pe] (e) {结束};
  \draw[op] (s) -- (p1);
  \draw[op] (p1) -- (p2);
  \draw[op] (p2) -- (p3);
  \draw[op] (p3) -- (p4);
  \draw[op] (p4) -- (p5);
  \draw[op] (p5) -- (p6);
  \draw[op] (p6) -- (p7);
  \draw[op] (p7) -- (p8);
  \draw[op] (p8) -- (pe);
  \draw[op] (pe) -- (e);
\end{tikzpicture}
}
\end{center}

\textbf{复杂度}：$O(V + E + F)$，每个元素遍历常数次。

\section{关键算法}

\subsection{流形边检查}

对每个面 $f$，遍历其边界顶点序列 $v_0 v_1 \dots v_{k-1} v_0$（$k=3$），
对每条无向边 $\{v_i, v_{i+1}\}$（端点排序作键）累加面计数：
\[
  \text{count}[\{a, b\}] \mathrel{+}= 1.
\]
最后扫描所有键，$\text{count} > 2$ 即为非流形边。

\begin{center}
\begin{tikzpicture}[
  vert/.style={circle, draw=blue!70!black, fill=blue!10, minimum size=6mm, inner sep=0pt, font=\footnotesize},
  he/.style={-, thick, draw=gray!60!black}
]
  % 中心边 a-b
  \node[vert] (a) at (0,0) {$a$};
  \node[vert] (b) at (2,0) {$b$};
  % 三个面共用此边（非流形）
  \node[vert] (c1) at (1,1.5) {$c_1$};
  \node[vert] (c2) at (1,-1.5) {$c_2$};
  \node[vert] (c3) at (3,0) {$c_3$};
  \draw[he] (a) -- (b);
  \draw[he] (a) -- (c1);
  \draw[he] (b) -- (c1);
  \draw[he] (a) -- (c2);
  \draw[he] (b) -- (c2);
  \draw[he] (a) -- (c3);
  \draw[he] (b) -- (c3);
  \node[below=1.8cm of a, xshift=1cm, font=\small, align=center, draw=red!60!black, fill=red!8]
    {无向边 $\{a, b\}$ 被 3 个面共享 $\Rightarrow$ \texttt{NonManifoldEdge}};
\end{tikzpicture}
\end{center}

\subsection{流形顶点检查}

对顶点 $v$，用 \texttt{VertexRing} 收集其所有 outgoing 半边。
对每条 outgoing 半边 $h$，判定是否为「边界半边」：
\[
  \text{isBoundary}(h) = \bigl(h.\mathrm{face} = \texttt{None}\bigr)
                     \vee \bigl(h.\mathrm{twin}.\mathrm{face} = \texttt{None}\bigr).
\]
统计 \texttt{boundary\_count}：
\begin{itemize}
  \item $= 0$：内部顶点，outgoing 环闭合 $\Rightarrow$ 合法；
  \item $= 2$：边界顶点，outgoing 环为开链，两端为边界半边 $\Rightarrow$ 合法；
  \item 其他：非流形（多条边界扇区或孤立边界）。
\end{itemize}

\begin{center}
\begin{tikzpicture}[
  vert/.style={circle, draw=blue!70!black, fill=blue!10, minimum size=6mm, inner sep=0pt, font=\footnotesize},
  he/.style={-Stealth, thick, draw=gray!60!black},
  bdhe/.style={-Stealth, thick, draw=red!70!black, dashed}
]
  % 左：内部顶点 boundary_count=0
  \node[vert] (v1) at (0,0) {$v$};
  \node[vert] (w1) at (90:1.6) {$w_1$};
  \node[vert] (w2) at (210:1.6) {$w_2$};
  \node[vert] (w3) at (330:1.6) {$w_3$};
  \draw[he] (v1) -- (w1);
  \draw[he] (v1) -- (w2);
  \draw[he] (v1) -- (w3);
  \draw[he] (w1) -- (w2);
  \draw[he] (w2) -- (w3);
  \draw[he] (w3) -- (w1);
  \node[below=2.4cm of v1, font=\small, align=center] {内部顶点\\$\text{bdry}=0$ ✓};

  % 右：边界顶点 boundary_count=2
  \node[vert] (v2) at (5,0) {$v$};
  \node[vert] (x1) at (5,1.6) {$x_1$};
  \node[vert] (x2) at (3.5,-0.8) {$x_2$};
  \node[vert] (x3) at (6.5,-0.8) {$x_3$};
  \draw[he] (v2) -- (x1);
  \draw[he] (v2) -- (x2);
  \draw[he] (v2) -- (x3);
  \draw[he] (x2) -- (x1);
  \draw[he] (x1) -- (x3);
  \draw[bdhe] (x2) -- (v2);
  \draw[bdhe] (x3) -- (v2);
  \node[below=2.4cm of v2, font=\small, align=center] {边界顶点\\$\text{bdry}=2$ ✓};
\end{tikzpicture}
\end{center}

\subsection{入口字段一致性}

\textbf{顶点入口}：$v.\mathrm{halfedge}$ 必须指向一条半边 $h$ 满足
$h.\mathrm{twin}.\mathrm{vertex} = v$（即 $h$ 的 origin 是 $v$）。

\textbf{面入口}：$f.\mathrm{halfedge}$ 必须指向一条半边 $h$ 满足
$h.\mathrm{face} = f$。

这两项检查能捕获 \texttt{split\_edge} / \texttt{flip\_edge} 等拓扑操作中
「忘记更新入口字段」的常见 bug。

\section{退化守卫}

\begin{itemize}
  \item \texttt{validate\_halfedge\_references}：若 \texttt{twin/next/prev} 为
        \texttt{None}，跳过该项检查（\texttt{None} 不视为悬空）；
  \item \texttt{validate\_face\_boundaries}：设 \texttt{max\_iter} =
        \texttt{halfedge\_count + 1}，防止 next 链成环时无限循环；
  \item \texttt{validate\_degenerate\_faces}：若面边界环非 3 顶点，跳过
        （已由 \texttt{validate\_face\_boundaries} 报告）。
\end{itemize}

\section{测试覆盖}

\begin{center}
\renewcommand{\arraystretch}{1.18}
\begin{tabular}{@{}ll@{}}
  \toprule
  \textbf{测试名} & \textbf{验证点} \\
  \midrule
  \texttt{clean\_triangle\_passes\_validation}       & 干净三角形 0 错误 \\
  \texttt{clean\_closed\_fan\_passes\_validation}    & 闭合扇形 0 错误（内部顶点） \\
  \texttt{detects\_twin\_mismatch}                   & 检测 twin 不互指 \\
  \texttt{detects\_dangling\_vertex\_reference}      & 检测悬空顶点引用 \\
  \texttt{detects\_degenerate\_face}                 & 检测共线三角形 \\
  \texttt{detects\_non\_triangular\_face}            & 检测四边形面 \\
  \texttt{detects\_vertex\_halfedge\_inconsistency}  & 检测入口 origin 不匹配 \\
  \texttt{topology\_ops\_preserve\_validity}         & split 后仍通过校验 \\
  \texttt{flip\_preserves\_full\_validity}           & flip 后仍通过校验 \\
  \texttt{validate\_first\_error\_passes\_on\_clean\_mesh}    & 干净网格返回 \texttt{None} \\
  \texttt{validate\_first\_error\_returns\_some\_on\_invalid\_mesh} & 错误网格返回首个违例 \\
  \bottomrule
\end{tabular}
\end{center}

\end{document}
